Large language models (LLMs) are increasingly being integrated into software engineering (SE) research and practice, yet their non-determinism, opaque training data, and evolving architectures complicate the reproduction and replication of empirical studies.
We present a community effort to scope this space, introducing a taxonomy of LLM-based study types together with eight guidelines for designing and reporting empirical studies involving LLMs.
The guidelines present essential (\must) criteria as well as desired (\should) criteria and target transparency throughout the research process.
Our recommendations, contextualized by our study types, are: (1) to declare LLM usage and role; (2) to report model versions, configurations, and fine-tuning; (3) to document tool architectures; (4) to disclose prompts and interaction logs; (5) to use human validation; (6) to employ an open LLM as a baseline; (7) to use suitable baselines, benchmarks, and metrics; and (8) to openly articulate limitations and mitigations.
Our goal is to enable reproducibility and replicability despite LLM-specific barriers to open science.
We maintain the study types and guidelines online as a living resource for the community to use and shape (\href{https://llm-guidelines.org/}{llm-guidelines.org}).
